\section{Adaptations to the PSO routing heuristic}
\label{app:pso}

Our drone routing builds on the PSO-inspired heuristic for the Team Orienteering
Problem (TOP) of Dang, Guibadj and Moukrim~\cite{PSO}. We summarise that algorithm
briefly and then describe the changes we made so that it runs fast enough for
large, grid-based wildfire instances. The original method does not represent each
drone route separately. Instead, a candidate solution is encoded as a single
permutation of all the locations to be considered, called a giant tour. A
\emph{split} procedure reads this permutation in order and, using dynamic
programming, cuts it into the set of feasible routes (one per drone) that collects
the most reward while respecting the battery budget. The algorithm then improves
solutions in the style of particle swarm optimization, recombining promising
permutations and refining them with local search moves that relocate a single
location (\emph{shift}) or exchange two of them (\emph{swap}). Each candidate move
is scored by re-running the split, so the split and the local search together
account for almost all of the running time.

We instantiate this framework on the wildfire grid: the locations are grid cells
carrying a wildfire-risk reward, the depots are charging stations, and travel cost
is the number of grid steps (Chebyshev distance) between cells. We also extend the
formulation so that within each battery cycle a drone may depart from and return to
any charging station rather than a single fixed depot. The giant tour therefore
contains several depot markers, and the split is free to assign each route to a
different station.

The remaining changes are accelerations that exploit the grid structure without
altering which solutions the search explores. First, we observe that a route can
only begin at a charging-station marker, and these markers are few compared with
the number of grid cells. The original split nonetheless scans every position of
the permutation; we instead run the dynamic program over depot positions only,
which lowers the cost of evaluating one candidate from roughly $\mathcal{O}(n\,m)$
to about $\mathcal{O}(k\,m \log k)$, where $n$ is the number of cells, $m$ the
number of drones, and $k$ the (small) number of depot markers. For typical
instances this alone speeds up evaluation by around sevenfold. Second, during local
search we derive simple conditions under which a candidate swap or shift cannot
possibly change the split outcome (for example, moves confined to cells that no
route currently visits), and we skip the expensive re-evaluation of such moves.
Depending on how densely the routes fill the grid, this discards between $40\%$ and
$85\%$ of candidate moves with no effect on the result. Third, rather than rebuild
every route after each accepted move, we cache the routes and recompute only the
ones a move actually touches, skipping the dynamic program entirely when nothing
relevant changes (about three quarters of the time); this makes the local search
roughly four times faster. Finally, we replace per-lookup dictionary access to
travel costs with a dense distance matrix and reuse pre-allocated buffers in the
inner loops, removing memory-allocation overhead from the hot path.

Battery feasibility is enforced inside the split itself: a route stops collecting
cells as soon as the drone could no longer return to a charging station within its
budget, so the procedure never produces an infeasible route. A greedy step then
extends any route that the search left short, allowing each drone to use its full
battery budget before returning. Each of the accelerations above was verified to
leave solution quality unchanged relative to a direct re-evaluation, so the reported
routing results reflect the same search as the baseline method, obtained at
substantially lower computational cost. The exact implementation of the PSO-based routing heuristic is available in the code (see the Code Availability statement).
